\subsection{Ungeeignete und eingeschränkt geeignete Projekttypen}
\label{subsec:ungeeignete-projekttypen}

Eingeschränkt geeignet sind Projekte mit umfangreicher Betriebssystemintegration,
plattformspezifischer UI, systemnahen Diensten oder Treibern. Die Tauri-Hülle
belegt dafür keine ausreichende native Basis. Auch stark verteilte
Infrastruktur sowie besondere Event-Streaming- oder HPC-Anforderungen können
eine andere Architektur erfordern; die Baseline bleibt höchstens für einen
abgegrenzten Teil nutzbar.

Hard-Real-Time-Systeme, grafikintensive Spiele, Firmware, Embedded-Systeme und
native Mobilanwendungen gehören nicht zum primären Zielbereich. Gewöhnliche
Live-Aktualisierungen machen eine Business-Anwendung nicht zum
Echtzeitsystem. Einzelne mobile oder hardwarenahe Funktionen sind nicht
ausgeschlossen; ihre technische Basis liefert das Template jedoch nicht.
Entscheidend ist daher nicht ein einzelnes Merkmal, sondern ob der technische
Kern des Produkts von der Baseline getragen wird.

Hohe Sicherheits- oder Compliance-Anforderungen schließen die Baseline nicht
aus, verlangen aber zusätzliche Evidenz für Schutzbedarf, Audits,
Zertifizierung und Rechtskonformität. Belegt sind Compose-Validierung und
Master-Images, nicht Compose-Start oder produktive Bereitstellung
\parencite{kleiveist2026ci-core-run}. Linux- und macOS-Pakete wurden
unsigniert gebaut; der Windows-Paketpfad ist fehlgeschlagen
\parencite{kleiveist2026ci-desktop-run}. Signing und Notarisierung bleiben
produktspezifisch \parencite{kleiveist2026release}. Auch die benötigten
TypeScript-, Python-, Rust-, Datenbank- und Deploymentkompetenzen sind zu
prüfen.
Fehlen diese Nachweise oder Kompetenzen, ist das Projekt höchstens bedingt
geeignet.
\beginnerinput{workspace/statement/800_einsatz_transfer/820}
